\paragraph{C.3u (Proved under the stated cone hypotheses): ordinary degree-one surjectivity in every characteristic.}
Let $k$ be an algebraically closed field of arbitrary
characteristic. Let $V$ be a smooth projective integral
$k$-variety of dimension $d\ge3$, and let $L$ be a
very ample invertible sheaf such that its full section
ring
\[
 S=\bigoplus_{a\ge0}H^0(V,L^a)
\]
is generated in degree one. Write $\mathfrak v=S_{>0}$,
and form the normal strictly henselian cone local ring
and its puncture
\[
 R=(S_{\mathfrak v})^{\mathrm{sh}},\qquad
 X=\operatorname{Spec}R\setminus\{\mathfrak m\},
 \qquad N=d+1\ge4.
\]
The geometric residue field in this strict
henselization is the given field $k$. Let
$0\ne f\in\mathfrak m$, and suppose
$Y=X\cap V(f)$ is regular. The element $f$ need
not be homogeneous. Let $\rho:X\to V$ be the
cone projection, and put
$A=(\operatorname{Pic}^0_{V/k})_{\mathrm{red}}$.
For ordinary algebraic equivalence defined by
actual line bundles on products with smooth
connected $k$-curves, the following statements hold:
\begin{align}
 A(k)&\xrightarrow[\rho^*]{\ \sim\ }
                    \operatorname{Pic}(X)_{\mathrm{alg}},
                    \label{eq:C3-cone-ordinary-source}\\
 \operatorname{Pic}(X)_{\mathrm{alg}}
       &\longrightarrow\operatorname{Pic}(Y)_{\mathrm{alg}}
                    \quad\text{is surjective}.
                    \label{eq:C3-cone-ordinary-surjectivity}
\end{align}
In characteristic zero the map in
\eqref{eq:C3-cone-ordinary-surjectivity} is an
isomorphism. In characteristic $p>0$ its kernel
is a finite $p$-primary group. Consequently
\[
                       J^1(X)\xrightarrow{\ \sim\ }J^1(Y)
\]
with rational coefficients in every characteristic.
These are statements about the specified smooth
cone source and arbitrary regular divisor punctures
inside it; they do not impose a cone structure
on the divisor or assert the same conclusion
for a deformation of the source cone.

\begin{proof}
\emph{1. The proper resolution and its actual
Picard family.}
The full section ring of the normal projective
variety is normal. Its degree-one generation
identifies the standard contraction
\[
 Z_0=\operatorname{Tot}(L^{-1})
          \longrightarrow\operatorname{Spec}S
\]
with the blow-up of the vertex. One can check
this on a degree-one homogeneous chart: the
blow-up chart is the affine-line bundle obtained
by retaining the radial coordinate over the
corresponding affine open of $\operatorname{Proj}S=V$.
Thus this contraction is proper, its exceptional
divisor is the zero section $V$, and it is an
isomorphism away from $V$. Base change gives
\[
 Z=Z_0\times_{\operatorname{Spec}S}\operatorname{Spec}R
                  \longrightarrow\operatorname{Spec}R.
\]
The scheme $Z$ is regular and integral. Indeed
it is the blow-up of the domain $R$, and its
regularity follows from that of $Z_0$ by the
localization and ind-etale base change defining
$R$. It is proper over the henselian ring $R$,
has the single prime exceptional divisor $V$,
and has reduced closed fibre exactly $V$.
Its complement is $X$, which is consequently
regular. The projection of the line bundle
induces $\pi:Z\to V$, extending $\rho$, with
\begin{equation}\label{eq:C3-cone-exceptional-restriction}
 \pi|_V=1_V,\qquad
             \mathcal O_Z(V)|_V=L^{-1}.
\end{equation}

The variety $A$ is an abelian variety. Moreover
there is a Poincare line bundle on $V\times A$,
normalized at a chosen $k$-point of $V$ and at
$0\in A$. Here the precise standard inputs are
\citet[Theorem 2.5(2), Exercise 4.3 with its answer,
and Theorem 5.4, author-version pp.~18, 27,
37--38, 65]{D_KleimanPicard2005v1}.
The section provided by the chosen point makes
the universal line bundle exist; the normality
of $V$ makes the reduced identity component
proper. Pulling this bundle back along $\rho$
gives an actual family on $X\times A$ and a
homomorphism
\[
           j:A\longrightarrow G,
           \qquad G=\mathbf{Pic}^{\mathrm{loc}}_{R/k}.
\]
The local Picard scheme is the separated group
scheme locally of finite type described in
C.3n, with $G(k)=\operatorname{Pic}(X)$.
Its closed reduced image $P=j(A)$ is a proper
connected algebraic subgroup, hence an abelian
variety, and $P(k)=j(A(k))$.

The map $j$ is injective on $k$-points.
In fact, if $Q\in A(k)$ has $\rho^*Q$ trivial,
the line bundle $\pi^*Q$ is trivial off $V$.
A rational trivialization therefore gives
\[
                         \pi^*Q\simeq\mathcal O_Z(bV)
                              \quad(b\in\mathbb Z).
\]
Restricting to $V$ gives $Q\simeq L^{-b}$.
An algebraically trivial line bundle has degree
zero on every curve, whereas a nonzero power
of the ample bundle $L$ does not. Thus $b=0$
and $Q$ is trivial. This argument asserts
injectivity on points; it does not assume
that $j$ has no infinitesimal kernel in
positive characteristic.

Every point of $A(k)$ gives an ordinary
algebraically trivial class on $X$. For example,
restrict the actual Poincare family to curves
joining the relevant points in the smooth
projective connected variety $A$, and pull
back to $X$. Such points can be joined by a
chain of integral curves, whose normalizations
are smooth projective over $k$. Conversely,
the equality between the points of
$\operatorname{Pic}^0_{V/k}$ and algebraically
trivial line bundles is
\citet[Proposition 5.10, pp.~41--42]{D_KleimanPicard2005v1}.

\emph{2. All source torsion at one suitable prime
comes from this actual family.}
The Neron--Severi group $\operatorname{NS}(V)$
is finitely generated in every characteristic;
see \citet[Theorem 7.3, p.~18]{D_ChengJiLarsonOlander2021}.
Choose a prime $\ell$ invertible in $k$ and not
dividing the order of the finite torsion subgroup
of
\[
                         \operatorname{NS}(V)/\mathbb Z[L].
\]
We claim that for every $n\ge1$,
\begin{equation}\label{eq:C3-cone-prime-torsion}
 A[\ell^n](k)\xrightarrow[\rho^*]{\ \sim\ }
                         \operatorname{Pic}(X)[\ell^n].
\end{equation}
Only surjectivity remains to be proved.

Let $M\in\operatorname{Pic}(X)[\ell^n]$.
The closure of the divisor of a rational
section extends $M$ to an invertible sheaf
$\widetilde M$ on the regular noetherian
integral scheme $Z$. The chosen triviality
of $M^{\ell^n}$ on $X$ gives
\[
                 \widetilde M^{\ell^n}\simeq\mathcal O_Z(bV)
                             \quad(b\in\mathbb Z).
\]
Put $Q=\widetilde M|_V$. By
\eqref{eq:C3-cone-exceptional-restriction},
$Q^{\ell^n}\simeq L^{-b}$. Hence the class
of $Q$ modulo $\mathbb Z[L]$ in
$\operatorname{NS}(V)$ is killed by $\ell^n$.
Our choice of $\ell$ makes that class zero.
Choose $a\in\mathbb Z$ with $[Q]=a[L]$,
and replace $\widetilde M$ by
$\widetilde M\otimes\mathcal O_Z(aV)$.
Its restriction, still denoted $Q$, now belongs
to $A(k)$. Its $\ell^n$th power is still of
the form $\mathcal O_Z(b'V)$.
Restricting this identity to $V$ shows
$\ell^n[Q]=-b'[L]$ in $\operatorname{NS}(V)$.
The left side is zero. Ampleness makes
$[L]$ of infinite order, so $b'=0$.
Consequently
\[
       \widetilde M^{\ell^n}\simeq\mathcal O_Z,
       \qquad Q^{\ell^n}\simeq\mathcal O_V.
\]

The bundle
$T=\widetilde M\otimes\pi^*Q^{-1}$ is
$\ell^n$-torsion and restricts trivially to
$V$. A trivialization of $T^{\ell^n}$ defines
a finite etale $\mu_{\ell^n}$-torsor over $Z$.
Its restriction to $V$ is trivial: after a
trivialization of $T|_V$, its defining unit
lies in $\Gamma(V,\mathcal O_V)^\times=k^\times$,
and has an $\ell^n$th root in $k$.
Now
\citet[Lemma 58.9.1, Tag 0A48]{D_StacksProperEtale}
gives an equivalence
\[
                  \operatorname{FEt}(Z)
                       \xrightarrow{\ \sim\ }
                  \operatorname{FEt}(V),
\]
because $Z$ is proper over the henselian
local ring $R$ and $V$ is its closed fibre.
Full faithfulness lifts a section of this
torsor over $V$ to a section over $Z$;
the torsor is therefore trivial. Hence
$T\simeq\mathcal O_Z$ and
$M\simeq\rho^*Q$, proving
\eqref{eq:C3-cone-prime-torsion}.
This uses the equivalence for finite etale
covers of the proper scheme $Z$, not a
claim that arbitrary line bundles or
parameter families algebraize from the
henselized product.

\emph{3. Ordinary source classes and target
surjectivity.}
We use the actual-family lemma proved in
C.3n: every ordinary algebraically trivial
class on the regular puncture lies in the
point group of a closed abelian subvariety
of its local Picard scheme. This lemma
starts with the given actual curve family
and constructs a dual-Jacobian family;
it does not infer an actual family from
an arbitrary subgroup of a local Picard
functor.

Apply it first to $X$, and let $B\subset G$
be such an abelian subvariety. Every point
of $B[\ell^n](k)$ is in $P(k)$ by
\eqref{eq:C3-cone-prime-torsion}.
The union of these torsion groups is Zariski
dense in $B$. Indeed its closed subgroup
closure, if of smaller dimension, could
have at most a fixed finite multiple of
$\ell^{2n\dim C}$ points of order dividing
$\ell^n$, where $C$ is its identity
component; this contradicts the
$\ell^{2n\dim B}$ such points of $B$.
Since $P$ is closed, $B\subset P$.
Together with the actual family in step 1
and pointwise injectivity of $j$, this
proves \eqref{eq:C3-cone-ordinary-source}.

For the target, write $\prod_a C_a$ for
the finite normalization of $R/(f)$ and
put
\[
 H=\prod_a\mathbf{Pic}^{\mathrm{loc}}_{C_a/k},
                  \qquad r:G\longrightarrow H.
\]
As in C.3n--p, the punctures of the
normal strictly henselian factors are the
disjoint components of $Y$. They have
dimension $N-1\ge3$, so $H$ is a
separated group scheme locally of finite
type and $H(k)=\operatorname{Pic}(Y)$.
The same actual-family lemma, applied to
the finitely many components, puts each
ordinary algebraically trivial class of
$Y$ in a closed abelian subvariety
$B\subset H$.

By C.3p, restriction gives isomorphisms
\[
 \operatorname{Pic}(X)[\ell^n]
       \xrightarrow{\ \sim\ }
 \operatorname{Pic}(Y)[\ell^n]
                      \qquad(n\ge1).
\]
Every point of $B[\ell^n](k)$ therefore
comes from $P(k)$ by step 2.
The image $r(P)$ is closed and proper,
since $P$ is proper and $H$ separated.
It contains the dense union of these
torsion groups, and hence contains $B$.
Every $k$-point of this image lifts to a
$k$-point of $P$: the nonempty fibre is
of finite type over the algebraically
closed field $k$. The points of $P$
are ordinary classes on $X$ by step 1.
This proves
\eqref{eq:C3-cone-ordinary-surjectivity}.

\emph{4. The integral kernel.}
The closed subgroup scheme
$K=\ker(r\circ j)\subset A$ is of
finite type. For every prime $\lambda$
invertible in $k$, C.3p and the pointwise
injectivity of $j$ imply
$K(k)[\lambda^n]=0$ for all $n\ge1$.
If $K$ had positive dimension, its reduced
identity component would be a
positive-dimensional abelian variety and
would have nonzero $\lambda$-torsion,
a contradiction. Thus $K$ is finite.
Its finite point group has no torsion
of order prime to the characteristic.
It is therefore zero in characteristic
zero and $p$-primary in characteristic
$p>0$. By
\eqref{eq:C3-cone-ordinary-source}, this
is exactly the ordinary restriction
kernel asserted above.
\end{proof}

The proof requires neither a splitting
of the whole local Picard identity
component nor the positive-characteristic
unipotent-kernel theorem used in C.3q.
The actual source family comes from
$\operatorname{Pic}^0(V)$, and proper
henselian finite-etale comparison detects
all the torsion needed for the density
argument. This closes the ordinary-family
issue for these cone sources in every
characteristic. C.3w addresses the more
general algebraizable characteristic-$p$
scope of C.3q by a separate alteration
and proper-curve construction.
